\section{Model and numerical transport solver}
\label{qt:computation}

\subsection{Kinetic model and transport problem}
Let $f(t,\mathbf x,\boldsymbol\xi)$ be the number density in position--velocity
space, with $n=\int f\,\mathrm d^3\xi$ and
$\mathbf u=n^{-1}\int\boldsymbol\xi f\,\mathrm d^3\xi$.
For normal, nonrelativistic gases with elastic binary collisions, the
Uehling--Uhlenbeck model \cite{qtUU} is
\begin{align}
 \partial_t f+\boldsymbol\xi\cdot\nabla f&=Q_\varepsilon(f),\notag\\
 Q_\varepsilon(f)&=\int B\bigl[
 f'f_1'(1+\varepsilon Af)(1+\varepsilon Af_1)
 -ff_1(1+\varepsilon Af')(1+\varepsilon Af_1')\bigr]
 \,\mathrm d^3\xi_1\,\mathrm d\Omega.
 \label{qt:uu}
\end{align}
Primes denote post-collision velocities; subscript 1 labels the second
particle. The signs $\varepsilon=+1,-1,0$ describe Bose, Fermi and classical
statistics. The constant $A=h_{\rm P}^3/(g_s m^3)$ converts number density
in velocity space to occupation per internal state, where $g_s$ is the
internal degeneracy and $h_{\rm P}$ is Planck's constant.
The collision measure is nonnegative and invariant under particle exchange
and time reversal. We assume isotropic scattering and sufficient angular
support to retain only the usual five collision invariants.

With $\qtU=\boldsymbol\xi-\mathbf u$, a regular local equilibrium is
\begin{equation}
 f^{(0)}=\frac1A\frac1{z^{-1}\exp[mU^2/(2k_{\mathrm B}T)]-\varepsilon},
 \qquad z=\exp[\mu_{\rm chem}/(k_{\mathrm B}T)].
\end{equation}
The classical form is understood as the $\varepsilon=0$ convention.
Collision balance holds at local equilibrium, while spatial gradients
generate a small correction. At first Chapman--Enskog order, rotational
symmetry reduces that correction to two unknown radial functions:
\begin{equation}
 \delta f=w\chi,\qquad
 \chi=-a(U)U_i\partial_i\ln T
       -b(U)D_{ij}(\qtU)\mathring S_{ij},\qquad
 D_{ij}=U_iU_j-\tfrac13U^2\delta_{ij}.
 \label{qt:ansatz}
\end{equation}
Here $w=f^{(0)}(1+A\varepsilon f^{(0)})$ and
$\mathring S_{ij}=(\partial_i u_j+\partial_j u_i)/2
-\delta_{ij}(\nabla\cdot\mathbf u)/3$ is the traceless strain rate.
Repeated Cartesian indices are summed. The functions $a$ and $b$ encode the response to a temperature
gradient and a shear, respectively. Symmetry determines their angular
structure, but it does not determine their radial dependence or their size.
The kinetic stress and heat flux are
\begin{equation}
 p_{ij}=m\int U_iU_j f\,\mathrm d^3\xi,
 \qquad q_i=\frac m2\int U^2U_i f\,\mathrm d^3\xi.
\end{equation}
Their first-order constitutive forms define the two coefficients:
$p_{ij}=p\delta_{ij}-2\mu\mathring S_{ij}$ and
$q_i=-\lambda\partial_iT$. At equilibrium, odd heat-flux moments vanish
and stress is isotropic; the correction supplies these dissipative fluxes.

Our numerical problem is therefore concrete: choose an equilibrium state and
a collision kernel, solve the two linear collision equations approximately,
and evaluate the moments that give $\lambda$ and $\mu$. These coefficients
are outputs, not fitted parameters. The calculation applies to dilute normal
gases with elastic binary collisions, close to local equilibrium. It does not
include a condensate, a superfluid component or a strongly interacting liquid.

Galerkin projection, Sonine--Laguerre expansions and Gaussian quadrature are
established numerical tools; the contribution is their implementation,
verification and application to the specified model. We do not implement
Wu's fast Fourier collision algorithm. We use a reference formula from that
work as one verification route \cite{qtWu}.

\subsection{Normalization: separating physics from units}
For particle mass $m$, the mass density is $\rho=mn$. We introduce
\begin{equation}
 v_0=\sqrt{\frac{k_{\mathrm B}T}{m}},\qquad
 \qtv=\frac{\qtU}{v_0},\qquad
 E=\frac{v^2}{2},\qquad
 \widehat f=\frac{v_0^3}{n}f.
 \label{qt:units}
\end{equation}
Thus $\int\widehat f\qtd^3v=1$. The normalization is conditional on the local
density: the physical $f$ integrates to $n$, while $\widehat f$ integrates
to one. All polynomial arguments below use the dimensionless energy $E$.

Define the statistical integrals by
\begin{equation}
 G_\nu^{(\varepsilon)}(z)
 =\frac1{\Gamma(\nu)}\int_0^\infty
 \frac{E^{\nu-1}}{z^{-1}e^E-\varepsilon}\qtd E.
 \label{qt:G}
\end{equation}
For $\varepsilon=0$, $G_\nu=z$. The integral definition also works for
Fermi fugacities greater than one, where a small-$z$ power series is unsuitable.
Writing
\begin{equation}
 D(E)=1-\varepsilon ze^{-E},\qquad
 \mathcal N=(2\pi)^{3/2}\frac{G_{3/2}^{(\varepsilon)}(z)}{z},
\end{equation}
the normalized equilibrium and its linearization weight are
\begin{equation}
 \widehat f_0=\frac{e^{-E}}{\mathcal N D(E)},\qquad
 \widehat A=z\mathcal N,\qquad
 r=1+\varepsilon\widehat A\widehat f_0=\frac1{D(E)},\qquad
 w=\widehat f_0r.
 \label{qt:equilibrium}
\end{equation}
In the rest of this numerical section $w$, the profiles and the collision
operator are dimensionless. The pressure is
\begin{equation}
 \qtZ=\frac{p}{nk_{\mathrm B}T}
 =\frac{G_{5/2}^{(\varepsilon)}(z)}{G_{3/2}^{(\varepsilon)}(z)}.
 \label{qt:Z}
\end{equation}
We use $0<z<1$ for normal Bose equilibria and $z>0$ for Fermi equilibria.
The classical limit is recovered as occupations become small. Equal fugacity
for two different statistics does not in general mean equal physical density
at fixed temperature: $n\lambda_{\mathrm{dB}}^3$ also depends on $G_{3/2}$
and on the internal-state degeneracy. The dimensionless comparisons below
isolate the transport calculation under the stated normalization.

\subsection{The collision kernel used for the first comparison}
Our baseline uses an isotropic hard-sphere-type differential cross section,
\begin{equation}
 B(g,\Omega)=g\frac{d^2}{4},\qquad
 g=|\qtU-\qtU_1|,\qquad
 \sigma_{\mathrm{tot}}=\int_{S^2}\frac{d^2}{4}\qtd\Omega=\pi d^2.
 \label{qt:kernel}
\end{equation}
Here $d$ is the hard-sphere diameter; the numerical convention sets $d=1$.
The kernel includes both the relative speed, which controls encounter rates,
and the differential cross section, which describes scattering into a solid
angle. Keeping exactly the same kernel when changing statistics isolates the
effect of the occupation factors. This is a model comparison, not a claim
that this cross section describes every atomic Bose or Fermi gas. In particular,
real fermionic scattering depends on spin composition and allowed partial waves.

The implementation also permits isotropic power laws $B=Cg^\gamma$,
$C>0$, $\gamma>-3$. The radial quadrature weight changes with $\gamma$;
the angular symmetry remains available. An arbitrary anisotropic scattering
law would require a further extension.

For the baseline convention the physical units can be restored using
\begin{equation}
 t_0=\frac1{nd^2v_0},\qquad \ell_0=\frac1{nd^2},\qquad
 \mu_{\mathrm{phys}}=\mu_{\mathrm{num}}\frac{\sqrt{mk_{\mathrm B}T}}{d^2},
 \qquad
 \lambda_{\mathrm{phys}}=\lambda_{\mathrm{num}}
 \frac{k_{\mathrm B}\sqrt{k_{\mathrm B}T/m}}{d^2}.
 \label{qt:restore}
\end{equation}
The scale $t_0$ is our unit of time; it is not an independently calculated
mean collision time, which also contains statistical and geometric factors.

\subsection{A quantum basis that satisfies matching}
With the equilibrium parameters held fixed, define the linearized operator by
\begin{equation}
 \qtL\chi=\left.\frac{\mathrm d}{\mathrm ds}
 Q_\varepsilon(f_0+s w\chi)\right|_{s=0},
\end{equation}
using the dimensionless variables and kernel in the numerical calculation.
The linearized equations for one representative vector component and one
off-diagonal tensor component are
\begin{equation}
 \qtL A_x=-wJ_x,\qquad J_x=v_x(E-\tfrac52\qtZ),\qquad
 \qtL B_{xy}=-wv_xv_y,
 \label{qt:profileeq}
\end{equation}
where $A_x=a(v)v_x$ and $B_{xy}=b(v)v_xv_y$.
There is no sum over $x$ or $y$ here. The tensor component suffices because
rotational symmetry makes every shear orientation equivalent. For a flow
$\partial_yu_x=G$, the two terms $\mathring S_{xy}=\mathring S_{yx}=G/2$
combine to give $b(v)v_xv_yG$, so no extra factor of two is missing.

We approximate the radial functions by $N$ profiles,
\begin{equation}
 a_N(v)=\sum_{n=1}^N\alpha_nP_n(E),\qquad
 b_N(v)=\sum_{n=1}^N\beta_nQ_n(E).
 \label{qt:expansion}
\end{equation}
For viscosity we take
\begin{equation}
 Q_n(E)=L_{n-1}^{(5/2)}(E).
\end{equation}
For heat conduction we explicitly remove the momentum component:
\begin{equation}
 \boxed{P_n(E)=L_n^{(3/2)}(E)-s_n,\qquad
 s_n=\frac{\int wv^2L_n^{(3/2)}(E)\qtd^3v}
                 {\int wv^2\qtd^3v}.}
 \label{qt:shift}
\end{equation}
Indeed, substituting this definition gives $\int wv^2P_n\qtd^3v=0$.
Therefore every finite linear combination $a_N$ satisfies thermal momentum
matching, regardless of the coefficients $\alpha_n$. Thermal mass and energy
matching follow from oddness. In the viscous sector the tensor structure
ensures all matching conditions by parity and tracelessness.

For reference, the generalized Laguerre polynomials are generated by
\begin{align}
 L_0^{(\alpha)}(E)&=1,\qquad L_1^{(\alpha)}(E)=\alpha+1-E,\\
 (n+1)L_{n+1}^{(\alpha)}(E)
 &=(2n+\alpha+1-E)L_n^{(\alpha)}(E)
 -(n+\alpha)L_{n-1}^{(\alpha)}(E).
\end{align}
In the classical case, $s_n=0$ for $n\geq1$ by Laguerre orthogonality.
In the quantum case the shifts are generally nonzero. For example,
the weighted moment identities yield
\begin{equation}
 s_1=\frac52(1-\qtZ),\qquad
 P_1(E)=\frac52\qtZ-E.
 \label{qt:firstprofile}
\end{equation}
Thus the first thermal vector is $P_1v_x=-J_x$. This is a useful explicit
example of how matching changes a classical polynomial in a quantum gas.

These shifted polynomials need not be mutually orthogonal under the quantum
weight. They only need to span an admissible, linearly independent trial space.
Galerkin projection handles their couplings through a full matrix. Orthogonality
of the classical polynomials would not diagonalize the collision operator
either. Increasing $N$ improves radial resolution; it does not increase the
Chapman--Enskog order, which remains first order in spatial gradients.

\subsection{From the integral equation to a matrix problem}
Define $\Phi_n^T=P_n(E)v_x$ and $\Phi_n^S=Q_n(E)v_xv_y$, and use the bracket
\begin{equation}
 [\phi,\psi]=-\int\phi\qtL\psi\qtd^3v
 =\frac14\int BW\,\Delta\phi\,\Delta\psi
          \qtd^3v\qtd^3v_1\qtd\Omega,
 \quad \Delta\phi=\phi+\phi_1-\phi'-\phi_1'.
 \label{qt:bracket}
\end{equation}
Here $W=\widehat f_0\widehat f_{0,1}r'r_1'$ is the detailed-balance weight
of the entire collision quadruple. It is distinct from the one-velocity
weight $w$ in the matching conditions.

For example, the thermal residual is
$R_N=\qtL(\sum_n\alpha_n\Phi_n^T)+wJ_x$.
Imposing $\int\Phi_m^TR_N\qtd^3v=0$ for each $m$ gives
\begin{equation}
 \boxed{M^T\boldsymbol\alpha=\boldsymbol d^T,\qquad
 M^T_{mn}=[\Phi_m^T,\Phi_n^T],\qquad
 d_m^T=\int w\Phi_m^TJ_x\qtd^3v.}
 \label{qt:galerkin}
\end{equation}
The superscript $T$ labels the thermal sector here, not a matrix transpose.
Similarly,
\begin{equation}
 M^S\boldsymbol\beta=\boldsymbol d^S,\qquad
 M^S_{mn}=[\Phi_m^S,\Phi_n^S],\qquad
 d_m^S=\int w\Phi_m^Sv_xv_y\qtd^3v.
\end{equation}
The bracket is nonnegative. Once the collision-invariant kernel is removed,
an admissible independent finite basis gives a positive-definite matrix under
the collision assumptions stated above. We solve these systems by Cholesky
factorization and check symmetry, positivity and residuals numerically.

The transport coefficients can then be evaluated in two equivalent ways:
\begin{equation}
 \boxed{\lambda_N=\sum_m d_m^T\alpha_m
 =\frac16\int wa_N(v)v^4\qtd^3v,\qquad
 \mu_N=\sum_m d_m^S\beta_m
 =\frac1{15}\int wb_N(v)v^4\qtd^3v.}
 \label{qt:coefficients}
\end{equation}
For the first equality for heat conduction, isotropy gives
$\int w a_Nv_x^2(E-5\qtZ/2)=\frac13\int wa_Nv^2(E-5\qtZ/2)$.
Matching removes the term proportional to $v^2$, and $E=v^2/2$ gives $1/6$.
For viscosity, $\int R(v)v_x^2v_y^2=\frac1{15}\int R(v)v^4$ for $x\ne y$.
This also explains why solving for a single component introduces no missing
factor of three.

For either sector, $d^{\mathsf T}M^{-1}d$ is the maximum over trial
coefficients $c$ of $2d^{\mathsf T}c-c^{\mathsf T}Mc$.
Nested trial spaces therefore give nondecreasing transport estimates when the
integrals are evaluated consistently and exactly. A small increment with $N$
is useful evidence of convergence, but it is not a rigorous error bound for
the infinite-dimensional solution.

\subsection{How the collision integrals are evaluated}
A direct bracket has eight integration variables: three for each incoming
velocity and two for the outgoing direction. Introduce centre and relative
velocities,
\begin{equation}
 \mathbf V=\frac{\qtv+\qtv_1}{2},\quad
 \mathbf q=\frac{\qtv-\qtv_1}{2},\quad
 \qtv=\mathbf V+\mathbf q,\quad \qtv_1=\mathbf V-\mathbf q.
\end{equation}
Elastic scattering replaces $\mathbf q$ by $\mathbf q'$ with $q'=q$.
The Jacobian is $\qtd^3v\qtd^3v_1=8\qtd^3V\qtd^3q$.
We can analytically average over the three angles of a common rotation,
because the equilibrium and kernel are isotropic.

Choose $\mathbf V$ along the polar axis and the incoming $\mathbf q$ in
a fixed meridian plane. Five variables remain:
\begin{equation}
 X=V^2,\quad Y=q^2,\quad
 a=\cos\theta_{\mathrm{in}},\quad b=\cos\theta_{\mathrm{out}},\quad
 \varphi\in[0,2\pi).
\end{equation}
The first two range over $[0,\infty)$, and $a,b\in[-1,1]$.
The incoming energies are
$E=(X+Y+2\sqrt{XY}a)/2$ and $E_1=(X+Y-2\sqrt{XY}a)/2$;
replace $a$ by $b$ for the outgoing energies. The azimuth still enters
the vector and tensor collision increments.

For $B=Cg^\gamma$, the reduced measure, including the factor $1/4$ in the
bracket and the normalization of the equilibrium, is
\begin{equation}
 4\,2^\gamma C\,
 X^{1/2}Y^{(1+\gamma)/2}e^{-X-Y}
 \mathcal R\qtd X\qtd Y\frac{\mathrm da}{2}\frac{\mathrm db}{2}
 \frac{\mathrm d\varphi}{2\pi},\qquad
 \mathcal R=\frac{1}{(G_{3/2}/z)^2DD_1D'D_1'}.
 \label{qt:measure}
\end{equation}
This formula shows precisely where the quantum modification is evaluated:
at all four collision energies, rather than only at one incoming velocity.
For the baseline $C=1/4$, $\gamma=1$, the numerical prefactor is $2$.

For vector profiles the rotational average of two selected-component
increments is $\Delta\mathbf A_m\cdot\Delta\mathbf A_n/3$.
For traceless symmetric tensor profiles the average for a selected
off-diagonal component is $\Delta\mathsf B_m:\Delta\mathsf B_n/10$.
These identities retain the component conventions used in the sources.

The two radial integrals use generalized Gauss--Laguerre rules with weights
$X^{1/2}e^{-X}$ and $Y^{(1+\gamma)/2}e^{-Y}$. The two polar cosines use
Gauss--Legendre rules; the periodic azimuth uses a trapezoidal rule. Gaussian
nodes and weights are constructed from the eigenproblem of the corresponding
Jacobi matrix. No sharp velocity cutoff is needed for these radial integrals.

A typical transport run uses $(n_r,n_p,n_\varphi)=(24,20,24)$, namely
$n_r^2n_p^2n_\varphi=5{,}529{,}600$ collision configurations. Near the Bose
threshold the factors $D(E)^{-1}$ vary more rapidly at low energy; the main
scan raises $n_r$ to $40$ for $z\geq0.8$. Separate refinement tests also
increase angular orders. This cost is paid to assemble small $N\times N$
matrices; solving the linear systems is not the expensive step.

\subsection{Verification and convergence are different questions}
\paragraph{First-profile verification.}
For a classical hard-sphere gas, the one-profile values in our units are
\begin{equation}
 \mu_1=\frac5{16\sqrt\pi},\qquad
 \lambda_1=\frac{75}{64\sqrt\pi}.
 \label{qt:classicalbenchmark}
\end{equation}
The quantum one-profile result can also be evaluated through the separate
integral representation in Eq.~(50) of Wu's work \cite{qtWu}, after matching
normalizations. Specifically, the quantum comparison uses ratios to the
classical one-profile values: constant scattering and spin multiplicities
in the original reference cancel in these ratios. We implement that reference
through a distinct, four-variable integration route. Across the classical
case, Bose $z=0.5,0.9$ and Fermi
$z=1,10$, the largest relative discrepancy in either coefficient is
$3.03\times10^{-5}$, or less than $0.0031\%$.
At $N=1$, Galerkin reduces to $M_{11}\alpha_1=d_1$ in each sector,
and the corresponding dimensionless coefficient is $d_1^2/M_{11}$.
Wu's variational ansatz has this same one-profile form; Eq.~(50) expresses
its result using reduced integrals $I_A$ and $I_B$. We evaluate those
integrals directly in a separate code path, rather than calling the bracket
routine. The two paths share equilibrium helper functions. This checks the
assembly against the same one-profile theory. It does not
establish $0.0031\%$ accuracy of a finite-$N$ solution to the full kinetic
equation, nor agreement with an experiment.

\paragraph{Basis and quadrature refinement.}
We separately vary $N$ at fixed integration settings, and refine the quadrature
at fixed $N=4$. Table~\ref{qt:convtable} reports the relative change from
$N=4$ to $N=5$, and the change between quadratures $(24,16,20)$ and
$(40,24,28)$. The latter uses $25{,}804{,}800$ collision configurations.
\begin{table}[htbp]\centering\small
\begin{tabular}{llrrrr}\toprule
 & & \multicolumn{2}{c}{$N=4\to5$ change [\%]}
 &\multicolumn{2}{c}{Quadrature change [\%]}\\
Statistics&$z$&$\mu$&$\lambda$&$\mu$&$\lambda$\\\midrule
\qtinput{\qtroot/data/convergence_table.tex}
\bottomrule\end{tabular}
\caption{Absolute relative changes, normalized to the coarser result.
Values printed as zero are below the displayed precision. These are convergence
diagnostics, not certified error bars.}\label{qt:convtable}
\end{table}
The largest displayed $N=4\to5$ change is about $0.3134\%$, in Fermi
conductivity at $z=10$. The basis truncation can therefore matter more than
the benchmark discrepancy. Figure~\ref{qt:convfig} uses $N=5$ as a common
reference to show the trend, without treating it as an exact solution.
\begin{figure}[htbp]\centering\qtfig{convergence}
\caption{Basis refinement for the classical gas, Bose $z=0.9$ and Fermi
$z=10$. Each curve is divided by its own $N=5$ result. The vanishing endpoint
is a consequence of this normalization, not proof of convergence at $N=5$.}
\label{qt:convfig}\end{figure}

\subsection{Transport results at equal approximation order}
Figure~\ref{qt:transportfig} compares each quantum result to the classical
result at the \emph{same} $N$. This avoids mixing a fourth-profile quantum
calculation with a first-profile classical denominator. The main scan contains
36 transport calculations: 18 statistical states evaluated at $N=1$ and $N=4$.
Representative $N=4$ values are listed in Table~\ref{qt:transporttable}.
\begin{figure}[htbp]\centering\qtfig{transport}
\caption{Quantum-to-classical transport ratios for the same isotropic kernel
and the same number of radial profiles in numerator and denominator.
Changing $N$ changes the numerical approximation, while changing $z$ changes
the equilibrium occupation. All displayed Bose states are above the
condensation threshold in the normal-gas description.}
\label{qt:transportfig}\end{figure}
\begin{table}[htbp]\centering\small
\begin{tabular}{llrrrr}\toprule
Statistics&$z$&$\mu_4$&$\lambda_4$&$\mu_4/\mu_{\mathrm{cl},4}$
 &$\lambda_4/\lambda_{\mathrm{cl},4}$\\\midrule
\qtinput{\qtroot/data/transport_table.tex}
\bottomrule\end{tabular}
\caption{Dimensionless coefficients and same-order ratios. Physical units
are restored by Eq.~\eqref{qt:restore}.}\label{qt:transporttable}
\end{table}
For this kernel and these states, Bose statistics lowers both coefficients
relative to the classical values, whereas Fermi statistics raises them.
The change in conductivity is appreciably smaller than the change in viscosity
over this scan. These are computed model results. Occupation factors alone
do not establish a universal monotonicity theorem for transport coefficients
under arbitrary changes of cross section and thermodynamic constraints.

\subsection{Implementation and reproducibility}
The solver is written in C++17 with Eigen for linear algebra and a JSON
library for structured results. It includes equilibrium moments, matched
polynomials, Gaussian rules, collision matrices, transport solves and the
dynamic extensions below. The numerical test executable reports 26 test
groups, supplemented by command-line smoke tests in CTest. Build commands,
local verification records and CI configuration accompany the source. These
software checks complement the physical benchmarks and convergence studies.

The migration from the earlier Python implementation was also checked against
saved results, retained as attributed regression fixtures. This is a
cross-implementation consistency check, not an independent physical benchmark.
Every figure in this note
is generated from the saved C++ outputs, without fitting or manually adjusting
the transport curves. The source files, calculation settings and tabulated
outputs are kept alongside the plotting data.
